\documentclass[12pt,a4paper]{article}
\usepackage[utf8]{inputenc}
\usepackage[T2A]{fontenc}
\usepackage[russian]{babel}
\usepackage{amsmath, amssymb, amsthm}
\usepackage{geometry}
\geometry{top=2cm, bottom=2cm, left=2.5cm, right=2.5cm}
\usepackage{graphicx}
\usepackage{hyperref}
\hypersetup{colorlinks=true, urlcolor=blue}

\newtheorem{theorem}{Теорема}
\newtheorem{definition}{Определение}
\newtheorem{criteria}{Критерий}

\title{Многоуровневая GRA-обнулёнка: статическая, циклическая и хаотическая динамика в задачах оптимизации и управления}
\author{Олег Битов}
\date{15 апреля 2026}

\begin{document}

\maketitle

\begin{abstract}
В работе представлено расширение концепции GRA-обнулёнки (Geometric Recursive Analytics) на три фундаментальных типа динамических систем: системы, сходящиеся к неподвижной точке (статическая GRA), системы с устойчивым предельным циклом (циклическая GRA) и системы с детерминированным хаосом (хаотическая GRA). Даны строгие математические критерии выбора типа обнулёнки на основе спектральных характеристик и показателей Ляпунова. Предложен гибридный подход для иерархических систем с разными временными масштабами. Показана связь с проектированием AGI/ASI, а также с культурными и философскими аспектами (инь-ян, дзен-гармония). Статья содержит полный математический аппарат, таблицы применимости и практические рекомендации.
\end{abstract}

\section{Введение}

GRA-обнулёнка изначально разрабатывалась как метод минимизации «пены» \(\Phi^{(l)}\) — меры отклонения многоуровневой системы от её целей \(G_l\). Классический подход предполагает сходимость к неподвижной точке, где \(\Phi^{(l)} \to 0\). Однако многие реальные системы (биологические, экономические, климатические) не могут быть полностью «успокоены»: они демонстрируют автоколебания или даже детерминированный хаос. В данной работе мы обобщаем GRA-обнулёнку на три типа динамических аттракторов и даём практические критерии выбора.

\section{Математическая основа GRA-обнулёнки}

\subsection{Иерархия уровней}

Рассмотрим систему, разбитую на уровни \(l = 0,1,\dots,K\). Каждый уровень характеризуется:

\begin{itemize}
    \item состоянием \(\Psi^{(l)} \in \mathcal{H}^{(l)}\);
    \item целью \(G_l\) — подпространством идеальных состояний;
    \item проектором \(\mathcal{P}_{G_l}\);
    \item пеной \(\Phi^{(l)} = \sum_{\mathbf{a}\neq\mathbf{b}} |\langle \Psi^{(\mathbf{a})}|\mathcal{P}_{G_l}|\Psi^{(\mathbf{b})}\rangle|^2\).
\end{itemize}

Общий функционал:
\[
J_{\text{total}} = \sum_{l=0}^{K} \Lambda_l \Phi^{(l)}, \quad \Lambda_l = \lambda_0 \alpha^l,\;0<\alpha<1.
\]

Цель GRA-обнулёнки — минимизация \(J_{\text{total}}\) путём изменения состояний \(\Psi^{(l)}\). Однако характер минимума зависит от динамики системы.

\subsection{Динамическая система и аттракторы}

Пусть эволюция системы задаётся ОДУ:
\[
\frac{d\mathbf{x}}{dt} = \mathbf{f}(\mathbf{x}, \mathbf{u}),
\]
где \(\mathbf{x}\) — вектор состояния (включает все \(\Psi^{(l)}\)), \(\mathbf{u}\) — управление. В отсутствие управления система стремится к одному из трёх типов аттракторов:

\begin{enumerate}
    \item \textbf{Неподвижная точка} \(\mathbf{x}^*\): \(\mathbf{f}(\mathbf{x}^*)=0\), асимптотически устойчива.
    \item \textbf{Предельный цикл} \(\Gamma\): периодическое решение \(\mathbf{x}(t+T)=\mathbf{x}(t)\).
    \item \textbf{Странный аттрактор} (детерминированный хаос): чувствительность к начальным условиям, непериодичность, ограниченность.
\end{enumerate}

Каждому типу соответствует свой вид GRA-обнулёнки.

\section{Статическая GRA (неподвижная точка)}

\subsection{Определение и математическая модель}

Статическая GRA применяется, когда целевым аттрактором является асимптотически устойчивая неподвижная точка. Пена определяется как квадратичное отклонение:
\[
\Phi_{\text{stat}} = \|\mathbf{x} - \mathbf{x}_{\text{goal}}\|^2.
\]
Управление (или эволюция) реализуется градиентным спуском:
\[
\dot{\mathbf{x}} = -\eta \nabla \Phi_{\text{stat}}.
\]

\subsection{Критерий применимости}

\begin{criteria}
Система может быть стабилизирована статической GRA, если в предполагаемой точке равновесия \(\mathbf{x}^*\) все собственные значения матрицы линеаризации \(A = \frac{\partial \mathbf{f}}{\partial \mathbf{x}}\big|_{\mathbf{x}^*}\) имеют отрицательные действительные части:
\[
\forall i:\ \operatorname{Re}(\lambda_i) < 0.
\]
\end{criteria}

\begin{theorem}[Сходимость статической GRA]
При выполнении указанного условия и достаточно малом \(\eta\) траектории сходятся к \(\mathbf{x}^*\) экспоненциально, и \(\Phi_{\text{stat}} \to 0\).
\end{theorem}

\subsection{Примеры}
\begin{itemize}
    \item Регулирование температуры.
    \item Стабилизация положения робота.
    \item Оптимизация логистических маршрутов.
\end{itemize}

\section{Циклическая GRA (предельный цикл)}

\subsection{Определение и математическая модель}

Циклическая GRA применяется, когда система должна совершать автоколебания (сердце, часы, ходьба). Цель — минимизировать амплитуду колебаний и стабилизировать частоту, но не подавлять их полностью. Введём фазу цикла \(\theta(t)\) и желаемую частоту \(\omega_0\). Пена:
\[
\Phi_{\text{cycle}} = \frac{1}{T}\int_0^T \|\mathbf{x}(t) - \mathbf{x}_{\text{ref}}(t)\|^2 dt + \mu (\dot{\theta} - \omega_0)^2.
\]
Управление:
\[
\dot{\mathbf{x}} = \mathbf{f}(\mathbf{x}) - \nabla \Phi_{\text{cycle}}.
\]

\subsection{Критерий применимости}

\begin{criteria}
Стационарная точка \(\mathbf{x}^*\) должна быть неустойчивой, но при бифуркации Андронова–Хопфа из неё рождается устойчивый предельный цикл. Условия:
\begin{enumerate}
    \item Матрица \(A\) имеет простую пару чисто мнимых собственных значений \(\lambda_{1,2}=\pm i\omega\).
    \item Все остальные \(\operatorname{Re}(\lambda_j)<0\).
    \item Первый ляпуновский коэффициент \(l_1 < 0\).
\end{enumerate}
\end{criteria}

При выполнении этих условий циклическая GRA может уменьшить амплитуду цикла, не уничтожая его.

\subsection{Примеры}
\begin{itemize}
    \item Стабилизация сердечного ритма.
    \item Управление шагающими роботами.
    \item Синхронизация генераторов в энергосети.
\end{itemize}

\section{Хаотическая GRA (странный аттрактор)}

\subsection{Определение и математическая модель}

Для хаотических систем (ураган, турбулентность) полное обнуление пены невозможно. Используется статистический подход: минимизируется средняя дисперсия или энтропия Колмогорова–Синая.
\[
\Phi_{\text{chaos}} = \langle \|\mathbf{x} - \langle \mathbf{x} \rangle \|^2 \rangle,
\]
где \(\langle \cdot \rangle\) — усреднение по времени. Управление может быть направлено на уменьшение \(\Phi_{\text{chaos}}\) или на подавление хаоса (перевод в режим цикла).

Метод управления хаосом (Пырагас):
\[
\dot{\mathbf{x}} = \mathbf{f}(\mathbf{x}) + K\bigl( \mathbf{x}(t-\tau) - \mathbf{x}(t) \bigr),
\]
где \(\tau\) — период желаемого цикла, \(K\) — матрица коэффициентов.

\subsection{Критерий применимости}

\begin{criteria}
Система хаотична, если её старший показатель Ляпунова \(\lambda_1 > 0\). Хаотическая GRA применима всегда, но её эффективность зависит от возможности управления параметрами.
\end{criteria}

\subsection{Примеры}
\begin{itemize}
    \item Прогноз и ослабление ураганов.
    \item Подавление турбулентности в трубах.
    \item Стабилизация хаотической сердечной аритмии.
\end{itemize}

\section{Гибридная GRA для иерархических систем}

В сложных системах (AGI/ASI, экономика, экосистемы) разные уровни иерархии могут требовать разных типов обнулёнки. Например:

\begin{table}[h]
\centering
\caption{Гибридный подход}
\begin{tabular}{|c|l|l|}
\hline
Уровень & Тип GRA & Пример \\
\hline
0 (быстрый) & Статическая & Стабилизация позы робота \\
1 (средний) & Циклическая & Шагающий цикл \\
2 (медленный) & Хаотическая & Адаптация к среде \\
\hline
\end{tabular}
\end{table}

Общий функционал:
\[
J_{\text{hybrid}} = \sum_{l=0}^{K} \Lambda_l \Phi_{\text{type}(l)}^{(l)}.
\]

\section{Связь с AGI/ASI и культурой}

\begin{itemize}
    \item \textbf{AGI/ASI}: рекурсивное самоулучшение требует баланса между эксплуатацией (статическая GRA) и исследованием (циклическая/хаотическая GRA). Идеальный AGI должен переключаться между режимами.
    \item \textbf{Культурная аналогия}: статика соответствует консерватизму, цикл — традиции и обновлению, хаос — революционным прорывам. Гармония (дзен) достигается при оптимальном соотношении \(80\%\) статики, \(20\%\) цикличности и малой доле хаоса.
\end{itemize}

\section{Заключение}

Разработана классификация видов GRA-обнулёнки, основанная на типах аттракторов динамической системы. Предложены строгие математические критерии для выбора статической, циклической, хаотической или гибридной GRA. Показано, что универсального «обнуления» пены не существует; в каждом случае нужно учитывать природу системы. Полученные результаты могут быть использованы при проектировании AGI/ASI, робототехнических систем, а также в экономике и биологии.

\section*{Благодарности}

Автор выражает благодарность сообществу GRA-исследователей за плодотворные обсуждения и вдохновение.

\begin{thebibliography}{9}
\bibitem{gra1} Многоуровневая архитектура GRA Мета-обнулёнка в мультиверсе. GitHub, 2026.
\bibitem{lyap} А.М. Ляпунов, «Общая задача об устойчивости движения», 1892.
\bibitem{hopf} E. Hopf, "Bifurcation of a periodic solution from a stationary solution", 1942.
\bibitem{pyragas} K. Pyragas, "Continuous control of chaos by self-controlling feedback", 1992.
\bibitem{alan} Аланский закон: этическая конституция для ИИ. GRA-Ethnos docs, 2026.
\end{thebibliography}

\end{document}